Here’s a clean, textbook-ready breakdown of what is well-established versus actively researched in zkML (with emphasis on R1CS / PLONK / Halo2 and ML circuits). I’ll keep it conceptual but precise.

1. What is firmly established

These are ideas that are stable, widely used, and unlikely to change fundamentally.

1.1 Arithmetic circuit view of ML

Neural networks can be compiled into algebraic constraints.

Linear layers → affine constraints
Activations → polynomial / lookup constraints
Entire model → composition of constraints

This reduction:

ML model
→
arithmetic circuit
→
proof system
ML model→arithmetic circuit→proof system

is foundational and settled.

1.2 R1CS as a canonical baseline

Even if not always used directly:

R1CS defines the minimal expressive core
Many compilers still lower to “R1CS-like” representations

👉 This is unlikely to disappear conceptually.

1.3 PLONKish paradigm (modern standard)

The shift to:

polynomial constraints
custom gates
lookup arguments

is now industry standard.

Key facts:

Lookup arguments are essential for efficiency
Custom gates reduce constraint count dramatically
Universal setup (vs circuit-specific) is a major advantage

👉 This design philosophy is mature and widely adopted.

1.4 Halo2-style circuit engineering

Well-established patterns:

columns (advice / fixed / selector)
regions and reusable gadgets
lookup-based nonlinearities
equality (copy) constraints

👉 This is now the practical way to build zk circuits.

1.5 Quantization is unavoidable

All zkML systems rely on:

𝑅
→
𝑍
→
𝐹
𝑝
R→Z→F
p
	​

fixed-point arithmetic
bounded integer ranges

👉 This is not a temporary workaround—it is structurally necessary.

1.6 Lookup-based nonlinearities

For functions like:

ReLU
sigmoid (quantized)
softmax (approximate)

Lookup tables are:

more efficient
more accurate (on discrete domain)

👉 This is now a dominant design pattern.

1.7 Constraint cost dominates performance

A fundamental law:

Proving cost ∝ number of constraints

Thus:

model size directly impacts proof time
architecture design affects feasibility

👉 This is a core engineering constraint.

2. What is still active research

These are areas where no consensus exists yet, and rapid progress is ongoing.

2.1 Scalable zkML for large models

Current reality:

MNIST / small CNN → feasible
ResNet / Transformers → very expensive

Open questions:

Can we prove inference for large foundation models efficiently?
How to reduce constraint count by orders of magnitude?
2.2 Better representations than circuits

Alternatives being explored:

zkVMs (prove execution instead of circuits)
STARK AIR systems
hybrid representations

Key question:

Is “circuit compilation” the right abstraction for ML?

2.3 Efficient nonlinearities

Still unresolved tradeoffs:

lookup vs polynomial vs hybrid
precision vs cost
domain size vs table size

Especially hard:

softmax
normalization layers
2.4 Training in zero knowledge

Inference is already hard—but training is much harder.

Challenges:

backpropagation doubles computation
gradients require high precision
massive memory requirements

👉 Fully practical zk training is still research-level.

2.5 Floating-point vs fixed-point gap

We currently prove:

quantized model
quantized model

But we want:

original floating-point model
original floating-point model

Open problems:

error guarantees
faithful equivalence proofs
stable quantization pipelines
2.6 Architecture design for zk

A new emerging field:

“zk-friendly neural networks”

Questions:

What activations are optimal?
What layer structures minimize constraints?
Can we co-design ML + proof system?
2.7 Proof system improvements

Even beyond circuits:

faster proving algorithms
better polynomial commitments
recursion and aggregation

Goal:

real-time or near real-time proofs
2.8 Data privacy vs model privacy

Different goals:

private input, public model
public input, private model
both private

Tradeoffs:

efficiency
expressiveness
cryptographic assumptions

Still evolving.

2.9 Verifiable training / auditing

Emerging direction:

prove dataset usage
prove fairness constraints
prove training correctness

This connects zkML to:

AI safety
governance
compliance
3. A clean conceptual split
Established core
ML
→
Arithmetic constraints
→
ZK proof
ML→Arithmetic constraints→ZK proof
	​

R1CS / PLONK / Halo2
quantization
lookup-based activations
Open frontier
How to scale, optimize, and generalize this pipeline
How to scale, optimize, and generalize this pipeline
	​

large models
better abstractions
training
architecture co-design
4. The most important meta-insight

There is a deep tension:

Machine learning prefers continuous, high-dimensional, approximate computation

Zero-knowledge proofs prefer discrete, algebraic, exact computation

Almost all current research is about:

bridging this gap efficiently
bridging this gap efficiently